\documentclass[11pt]{article}
\usepackage[margin=1in]{geometry}
\usepackage[T1]{fontenc}
\usepackage{graphicx}
\usepackage{hyperref}
\usepackage{enumitem}
\setlength{\parskip}{0.5em}
\setlength{\parindent}{0pt}

\title{diffusion models for video generation}
\author{arXiv GraphRAG --- phase-6 deep research}
\date{\today}

\begin{document}
\maketitle

\begin{abstract}
The review provides a solid historical foundation and covers many 2023-2024 papers, but misses the critical recent shift from U-Net to Transformer-based diffusion backbones and the role of LLMs in video orchestration.
\end{abstract}

\section{From Temporal Sequences to Generative Diffusion: A Survey of Video Synthesis Evolution}

The ability to synthesize realistic, high-fidelity video remains one of the most significant challenges in computer vision and generative artificial intelligence. Unlike static image synthesis, video generation requires the simultaneous modeling of complex spatial structures and intricate temporal dynamics, necessitating a robust understanding of the foundational principles within the Foundations of Generative Modeling and Video Representation. As the demand for automated content creation grows across digital media, robotics, and virtual simulations, the ability to produce motion that is both visually convincing and semantically coherent has become a critical driver of technological innovation.

This review traces the technological trajectory of temporal synthesis, beginning with the historical limitations of Early Video Generation: GANs and Sequence Models, which often struggled with temporal instability and limited diversity. We then examine the transformative paradigm shift brought about by The Rise of Video Diffusion Models, which leveraged iterative denoising processes to achieve unprecedented levels of detail. The discussion evolves to address the essential requirements for Controllability and Motion-Aware Synthesis, ensuring that generated movements align with specific user intents. Furthermore, we explore the ongoing efforts regarding Scaling, Efficiency, and Long-Term Consistency to overcome the computational bottlenecks inherent in high-resolution temporal modeling. Finally, the review concludes by examining The Current Frontier: Transformers and Advanced Video Diffusion, where the integration of large-scale architectural innovations is redefining the boundaries of generative video synthesis.

\subsection{Foundations of Generative Modeling and Video Representation}

The evolution of generative modeling has transitioned from early variational frameworks to sophisticated diffusion-based architectures. Initial progress was driven by the need to handle intractable posterior distributions in directed probabilistic models with continuous latent variables, a challenge addressed by the introduction of stochastic variational inference and the reparameterization of the variational lower bound in \texttt{arXiv:1312.6114}. This foundation enabled the optimization of complex models via standard stochastic gradient methods.

Parallel to variational approaches, Generative Adversarial Networks (GANs) introduced a minimax two-player game between a generator and a discriminator (\texttt{arXiv:1406.2661}). While GANs offered a different paradigm, they faced significant stability issues, specifically a lack of proven convergence during training. This was partially mitigated by the Two Time-Scale Update Rule (TTUR), which employs individual learning rates for the generator and discriminator to facilitate convergence to a local Nash equilibrium (\texttt{arXiv:1706.08500}). Further refinements in GAN architectures, such as the style-based generator (\texttt{arXiv:1812.04948}), improved image synthesis by enabling the unsupervised separation of high-level attributes and stochastic variation through latent code manipulation.

A significant shift occurred with the emergence of score-based modeling. To address the problem of ill-defined gradients when data resides on low-dimensional manifolds, \texttt{arXiv:1907.05600} proposed estimating gradients of the data distribution via score matching on noise-perturbed data, utilizing annealed Langevin dynamics. This approach was further unified with diffusion probabilistic models (DPMs) in \texttt{arXiv:2006.11239}, which connected nonequilibrium thermodynamics to a weighted variational bound and denoising score matching. While DPMs proved highly effective for image synthesis, they were hindered by slow sampling speeds due to the requirement of simulating a Markov chain over many steps. This limitation was directly addressed by \texttt{arXiv:2010.02502}, which introduced non-Markovian diffusion processes (Denoising Diffusion Implicit Models) to achieve sampling speeds 10x to 50x faster than standard DPMs using the same training procedure.

Modern generative capabilities have been expanded through multi-modal conditioning and latent space operations. The ability to bridge vision and language was advanced by \texttt{arXiv:2103.00020}, which utilized large-scale (image, text) pairs to enable zero-shot transfer, overcoming the limitations of fixed object categories. To combat the high computational costs of operating directly in pixel space, \texttt{arXiv:2112.10752} proposed applying diffusion models within the latent space of pretrained autoencoders. By incorporating cross-attention layers, this method allows for flexible conditioning on text or bounding boxes, facilitating tasks ranging from semantic scene synthesis to super-resolution.

\subsection{Early Video Generation: GANs and Sequence Models}

Before the dominance of diffusion-based architectures, early research in video generation focused on two distinct trajectories: sequence modeling for video-to-text tasks and Generative Adversarial Networks (GANs) for visual synthesis. Initial efforts in video understanding utilized recurrent structures to bridge the gap between visual content and language. For instance, \texttt{arXiv:1412.4729} addressed the scarcity of captioned video datasets by leveraging a unified deep neural network that transferred knowledge from over 1.2M category-labeled images to enable open-domain video description. While effective, subsequent work identified that local word generation often failed to capture holistic semantics. To address this, \texttt{arXiv:1505.01861} introduced the LSTM-E framework, which jointly modeled visual-semantic embeddings and sentence generation to ensure better alignment between subjects, verbs, and objects.

Parallel to these linguistic models, the development of GANs for video synthesis focused on untangling temporal dynamics from spatial content. Early architectures like \texttt{arXiv:1609.02612} utilized spatio-temporal convolutional structures to separate foreground from background for scene dynamics recognition. This concept of decoupling was further refined in the Temporal Generative Adversarial Nets (TGAN) proposed in \texttt{arXiv:1611.06624}, which employed a dual-generator architecture (temporal and image) and singular value clipping to mitigate the inherent instability of training complex GANs. Building on these foundations, \texttt{arXiv:1707.07493} (MoCoGAN) achieved unsupervised decomposition of motion and content using both image and video discriminators.

As the field progressed, the primary challenge shifted from basic generation to the computational scaling required for high-resolution and long-duration videos. \texttt{arXiv:1811.09245} addressed the exponential cost of resolution by using a stack of small sub-generators, allowing computational costs to scale linearly. Similarly, \texttt{arXiv:1907.06571} introduced the Dual Video Discriminator GAN (DVD-GAN), which utilized a computationally efficient discriminator decomposition to handle complex datasets.

A significant shift occurred when researchers moved away from discrete 3D RGB grids toward continuous representations. While \texttt{arXiv:2112.14683} (StyleGAN-V) achieved continuous-time synthesis through positional embeddings and holistic temporal aggregation, \texttt{arXiv:2202.10571} (DIGAN) addressed the limitations of grid-based scaling by utilizing Implicit Neural Representations (INRs). By manipulating differential space and time coordinates, DIGAN provided a more robust framework for modeling the spatio-temporal complexity and continuity required for long video generation. This evolution demonstrates a clear progression from simple recurrent sequence modeling to sophisticated adversarial frameworks capable of decomposing motion, scaling resolution, and eventually representing video as a continuous temporal signal.

\subsection{The Rise of Video Diffusion Models}

The evolution of video generation has been fundamentally shaped by the transition from text-to-image diffusion frameworks to specialized spatio-temporal architectures. Early foundations were laid by GLIDE, which demonstrated the efficacy of text-conditional diffusion using CLIP and classifier-free guidance for image synthesis and inpainting (\texttt{arXiv:2112.10741}). Building upon these image-centric principles, researchers sought to extend diffusion to the temporal dimension. Video Diffusion Models (VDM) established a baseline by extending standard image architectures through conditional sampling and joint training on both image and video data to mitigate gradient variance (\texttt{arXiv:2204.03458}).

A significant challenge in this transition is the scarcity of paired text-video datasets and the immense computational costs associated with training large-scale models (\texttt{arXiv:2205.15868}). To circumvent the need for massive text-video datasets, Make-A-Video proposed decomposing temporal U-Net and attention tensors into spatial and temporal approximations, allowing the model to leverage paired text-image data alongside unsupervised video footage (\texttt{arXiv:2209.14792}). In contrast, CogVideo addressed the scale problem by utilizing a 9B-parameter transformer and a multi-frame-rate hierarchical training strategy, inheriting knowledge from a pretrained text-to-image model (\texttt{arXiv:2205.15868}). While Make-A-Video focuses on data efficiency, Imagen Video pursues high-definition output through a cascade of fully-convolutional spatial and temporal super-resolution models (\texttt{arXiv:2210.02303}).

Efficiency and temporal consistency remain central themes in recent developments. MagicVideo and Latent Video Diffusion Models (LVDM) both utilize low-dimensional latent spaces to reduce computational demands (\texttt{arXiv:2211.13221}; \texttt{arXiv:2211.11018}). However, they differ in their architectural approach: MagicVideo employs a 3D U-Net with a frame-wise lightweight adaptor and a directed temporal attention module (\texttt{arXiv:2211.11018}), whereas LVDM introduces hierarchical diffusion in a 3D latent space to support long-form video generation of over 1000 frames, using conditional latent perturbation to combat error accumulation (\texttt{arXiv:2211.13221}).

Recent research has also pivoted toward zero-shot capabilities and structural control. Text2Video-Zero bypasses the need for additional training entirely by reprogramming frame-level self-attention in existing text-to-image models (\texttt{arXiv:2303.13439}). For more controlled synthesis, Structure and Content-Guided Video Synthesis achieves disentanglement by training on monocular depth estimates (\texttt{arXiv:2302.03011}). Finally, large-scale implementations like ModelScope integrate spatio-temporal blocks within a 1.7 billion parameter denoising U-Net to ensure smooth movement and consistent frame generation (\texttt{arXiv:2308.06571}).

\subsection{Controllability and Motion-Aware Synthesis}

Achieving high-fidelity video generation requires addressing the fundamental tension between maintaining spatial content integrity and introducing dynamic motion. Early advancements in spatial controllability were established by ControlNet, which utilizes zero-initialized convolution layers to inject spatial conditioning—such as depth or human pose—into pretrained text-to-image models without degrading existing features \texttt{arXiv:2302.05543}. This architectural principle of adding control without destroying learned weights serves as a foundation for several video-centric methodologies.

A primary challenge in video synthesis is the "intertwining" of spatial content and temporal dynamics, which increases modeling complexity \texttt{arXiv:2312.04483}. To mitigate this, HiGen proposes a hierarchical spatio-temporal decoupling paradigm, separating spatial reasoning from temporal reasoning through a unified denoiser \texttt{arXiv:2312.04483}. Similarly, AnimateDiff addresses the difficulty of adding motion to personalized text-to-image models by introducing a plug-and-play motion module and MotionLoRA, which learns transferable motion priors without requiring model-specific tuning \texttt{arXiv:2307.04725}. While AnimateDiff focuses on animating existing models, other approaches focus on extending latent diffusion models (LDMs) directly. For instance, Align your Latents extends pretrained image LDMs by introducing temporal layers to enable high-resolution synthesis \texttt{arXiv:2304.08818}, a strategy echoed by Stable Video Diffusion, which employs a three-stage training process—comprising text-to-image pretraining, video pretraining, and high-quality finetuning—to scale latent video diffusion to large datasets \texttt{arXiv:2311.15127}.

Beyond structural extensions, recent research emphasizes explicit motion guidance to improve temporal consistency. MoVideo addresses the failure of existing models to explicitly consider motion by incorporating depth and optical flow guidance in the latent space to regulate cross-frame correspondences \texttt{arXiv:2311.11325}. VideoComposer advances this compositional control by utilizing a Spatio-Temporal Condition (STC) encoder, which unifies spatial conditions (like sketches) with temporal conditions, such as motion vectors derived from compressed videos \texttt{arXiv:2306.02018}. This contrasts with FateZero, which focuses on zero-shot video editing rather than synthesis from scratch; FateZero utilizes attention fusion—blending self-attention and spatial-temporal attention maps—to ensure frame consistency during text-based editing without per-prompt training \texttt{arXiv:2303.09535}.

While these methods significantly improve controllability, they face varying degrees of complexity. For example, while the hierarchical decoupling in HiGen simplifies the reasoning process \texttt{arXiv:2312.04483}, the scaling strategies seen in Stable Video Diffusion require massive, systematically curated datasets to overcome the lack of unified training methods in the field \texttt{arXiv:2311.15127}. Collectively, the literature shows a transition from simple spatial conditioning \texttt{arXiv:2302.05543} toward complex, decoupled, and motion-aware architectures that treat temporal dynamics as a distinct, controllable dimension \texttt{arXiv:2306.02018}, \texttt{arXiv:2312.04483}.

\subsection{Scaling, Efficiency, and Long-Term Consistency}

The advancement of video diffusion models is currently constrained by high computational demands, data scarcity, and the difficulty of maintaining temporal coherence over extended durations. Recent research has addressed these challenges through architectural innovations, training-free refinement, and novel temporal modeling techniques.

Efficiency in video generation is primarily hindered by the quadratic complexity of attention layers and the high dimensionality of video data. To mitigate these costs, \texttt{arXiv:2403.14148} proposes a Content-Frame Motion-Latent Decomposition (CMD) approach, which utilizes a compact latent space to separate content frames from low-dimensional motion latents, thereby allowing the integration of pretrained image diffusion models. Similarly, \texttt{arXiv:2403.07711} seeks to overcome the quadratic scaling of attention by replacing it with linear-time bidirectional Structured State Spaces (SSMs), such as Mamba, to facilitate more efficient temporal feature extraction. Other frameworks attempt to reduce complexity by decoupling tasks; for instance, the ConFiner framework (\texttt{arXiv:2408.13423}) decomposes generation into structure control and spatial-temporal refinement using a chain of off-the-shelf diffusion experts.

Long-term consistency and temporal modeling present a separate set of challenges, particularly regarding the rigid sequence lengths of synchronous models and the error accumulation in asynchronous auto-regressive approaches. \texttt{arXiv:2503.07418} introduces AR-Diffusion, which employs a non-decreasing corruption timestep constraint and specialized schedulers (FoPP and AD) to bridge the training-inference gap, enabling flexible video length generation. For scenarios where training is computationally prohibitive, \texttt{arXiv:2503.09926} introduces VideoMerge, a training-free method that merges short videos from pretrained models to achieve extended durations. However, achieving smoothness and long-range consistency in such merged sequences remains a significant hurdle.

The precision of temporal modeling is further refined by addressing the limitations of scalar timestep variables. \texttt{arXiv:2410.03160} introduces a Vectorized Timestep Approach (VTV), allowing each frame to follow an independent noise schedule to capture fine-grained temporal dependencies, which helps prevent catastrophic forgetting during fine-tuning. In the realm of training-free guidance, \texttt{arXiv:2501.07563} utilizes a motion consistency loss to capture inter-frame feature correlation patterns, providing gradient-based guidance to improve temporal coherence.

Finally, the progress in these areas must be measured against standardized metrics. While existing metrics often fail to align with human perception, the VBench suite (\texttt{arXiv:2311.17982}) provides a hierarchical, 16-dimension benchmark to dissect video quality across disentangled dimensions. This is critical as models attempt to overcome data limitations, such as the high-quality data scarcity addressed by VideoCrafter2 (\texttt{arXiv:2401.09047}), which finetunes spatial modules with high-quality images to shift distributions without degrading motion.

\subsection{The Current Frontier: Transformers and Advanced Video Diffusion}

Recent advancements in video diffusion are increasingly defined by the transition toward Diffusion Transformers (DiTs) and the pursuit of efficient spatio-temporal modeling. A primary challenge in this transition is the entanglement of content and style features within DiT internal representations, which complicates tasks like style transfer [\texttt{arXiv:2601.06202}, \texttt{arXiv:2601.20175}]. To address this, researchers have proposed curriculum continual learning frameworks—training on mixtures of clean and noisy triplets to enable generalization to unseen styles without degrading content [\texttt{arXiv:2601.06202}, \texttt{arXiv:2601.20175}]. Building on this, VISTA introduces a specialized in-context framework using a lightweight style adapter and the VISTA-1000 synthetic dataset to mitigate temporal inconsistency and flickering artifacts often caused by long-term motion or occlusions [\texttt{arXiv:2605.17312}].

The architectural evolution of these models focuses heavily on the trade-off between computational cost and modeling depth. While Full 3D Attention provides comprehensive spatio-temporal modeling, it is prohibitively expensive, whereas Local Factorized Attention often fails to capture complex dynamics [\texttt{arXiv:2603.09721}]. FrameDiT attempts to bridge this gap by introducing Matrix Attention, which treats entire frames as matrices to enable efficient global spatio-temporal modeling [\texttt{arXiv:2603.09721}]. This drive for efficiency is echoed in broader surveys that categorize advancements into four paradigms: step distillation, efficient attention, model compression, and cache/trajectory optimization [\texttt{arXiv:2604.15911}].

Beyond architectural efficiency, recent methodologies focus on versatile control and high-fidelity synthesis. Temporal In-Context Fine-Tuning (TIC-FT) offers a way to achieve efficient adaptation for image-to-video and video-to-video tasks without requiring architectural changes, instead utilizing buffer frames of increasing noise levels to ensure smooth transitions [\texttt{arXiv:2506.00996}]. For temporal resolution, DiffuseSlide provides a training-free approach for high frame rate (FPS) generation by employing sliding window latent denoising and noise re-injection, specifically targeting the flickering and degradation common in long sequences [\texttt{arXiv:2506.01454}]. Furthermore, compositional synthesis has been advanced through object-centric learning, where pose-invariant slots are conditioned on pretrained models to allow for semantic editing, such as the insertion or deletion of objects [\texttt{arXiv:2507.20855}].

Collectively, these works highlight a shift from basic video generation toward highly controllable, efficient, and compositionally aware systems. However, significant hurdles remain, including the need for robust real-time long-horizon generation, the preservation of identity and motion coherence in image-to-video tasks [\texttt{arXiv:2605.17248}], and the ongoing struggle to balance computational efficiency with high-quality visual understanding [\texttt{arXiv:2502.07001}, \texttt{arXiv:2504.16081}].

\subsection{Tensions and Open Problems}

A primary tension in the literature concerns the architectural foundation of video synthesis: whether models should prioritize global temporal coherence through unified diffusion or achieve high-fidelity motion through decoupled, object-centric control. Early large-scale approaches, such as those presented in \texttt{arXiv:2210.02303}, focus on high-definition diffusion to maintain visual quality, yet they often struggle with fine-grained compositional control. In contrast, \texttt{arXiv:2507.20855} argues that such monolithic frameworks lack the necessary granularity for complex scenes, proposing instead a temporal object-centric learning approach to manage individual entity dynamics. This debate extends to the trade-off between simplicity and controllability; while \texttt{arXiv:2209.14792} demonstrates that text-to-video can be achieved without massive text-video datasets, more recent critiques in \texttt{arXiv:2302.07685} suggest that such methods may fail to capture the probabilistic complexities of latent space, necessitating more structured projection techniques to ensure temporal stability.

Another significant debate centers on the evaluation and benchmarking of generative quality. For years, the field relied on metrics derived from early adversarial frameworks, such as the convergence properties discussed in \texttt{arXiv:1706.08500}. However, recent scholarship suggests these metrics are insufficient for the nuances of video. Specifically, \texttt{arXiv:2311.17982} critiques the reliance on older video metrics like those in \texttt{arXiv:1812.01717}, arguing that they fail to capture the holistic semantic and temporal dimensions required for modern generative models. This tension highlights a growing gap between the rapid advancement of model capabilities and the stagnant nature of the metrics used to validate them, creating an open problem regarding how to objectively quantify "motion realism" versus "visual fidelity."

Finally, there is an ongoing disagreement regarding the optimal path toward long-form video generation: whether to scale existing latent diffusion models or to move toward asynchronous, auto-regressive architectures. While models like those in \texttt{arXiv:2311.15127} attempt to scale latent video diffusion to achieve stability, they are often criticized for computational inefficiency and a lack of structural decoupling, as noted in \texttt{arXiv:2408.13423}. This has led to a bifurcation in research, where some favor the massive scaling of unified transformers, while others advocate for the asynchronous, modular approaches seen in \texttt{arXiv:2503.07418}, which attempt to break the temporal dependency bottleneck.

\subsection{Outlook}

The frontier of video generation is rapidly shifting from mere pixel synthesis toward sophisticated, content-preserving style transfer and efficient, high-frame-rate generation. As the field moves beyond the limitations of standard diffusion, the integration of Diffusion Transformers (DiT) with advanced attention mechanisms, such as the matrix attention explored in \texttt{arXiv:2603.09721}, promises to resolve the long-standing conflict between computational efficiency and temporal resolution. The next generation of models will likely move toward "asynchronous intelligence," where video is not generated as a single block, but as a continuous, evolving stream of temporally consistent and stylistically controllable representations.

\section*{References}
\begin{enumerate}\small
  \item UCF101: A Dataset of 101 Human Actions Classes From Videos in The Wild (2012). \texttt{arXiv:1212.0402}
  \item Auto-Encoding Variational Bayes (2013). \texttt{arXiv:1312.6114}
  \item Translating Videos to Natural Language Using Deep Recurrent Neural
  Networks (2014). \texttt{arXiv:1412.4729}
  \item Coherent Multi-Sentence Video Description with Variable Level of Detail (2014). \texttt{arXiv:1403.6173}
  \item Generative Adversarial Networks (2014). \texttt{arXiv:1406.2661}
  \item Deep Unsupervised Learning using Nonequilibrium Thermodynamics (2015). \texttt{arXiv:1503.03585}
  \item Sequence to Sequence -- Video to Text (2015). \texttt{arXiv:1505.00487}
  \item U-Net: Convolutional Networks for Biomedical Image Segmentation (2015). \texttt{arXiv:1505.04597}
  \item Using Descriptive Video Services to Create a Large Data Source for Video
  Annotation Research (2015). \texttt{arXiv:1503.01070}
  \item Jointly Modeling Embedding and Translation to Bridge Video and Language (2015). \texttt{arXiv:1505.01861}
  \item Video Paragraph Captioning Using Hierarchical Recurrent Neural Networks (2015). \texttt{arXiv:1510.07712}
  \item The Long-Short Story of Movie Description (2015). \texttt{arXiv:1506.01698}
  \item Hierarchical Recurrent Neural Encoder for Video Representation with
  Application to Captioning (2015). \texttt{arXiv:1511.03476}
  \item MovieQA: Understanding Stories in Movies through Question-Answering (2015). \texttt{arXiv:1512.02902}
  \item Generating Videos with Scene Dynamics (2016). \texttt{arXiv:1609.02612}
  \item Temporal Generative Adversarial Nets with Singular Value Clipping (2016). \texttt{arXiv:1611.06624}
  \item TGIF: A New Dataset and Benchmark on Animated GIF Description (2016). \texttt{arXiv:1604.02748}
  \item Attention Is All You Need (2017). \texttt{arXiv:1706.03762}
  \item GANs Trained by a Two Time-Scale Update Rule Converge to a Local Nash
  Equilibrium (2017). \texttt{arXiv:1706.08500}
  \item Decoupled Weight Decay Regularization (2017). \texttt{arXiv:1711.05101}
  \item Neural Discrete Representation Learning (2017). \texttt{arXiv:1711.00937}
  \item MoCoGAN: Decomposing Motion and Content for Video Generation (2017). \texttt{arXiv:1707.04993}
  \item Quo Vadis, Action Recognition? A New Model and the Kinetics Dataset (2017). \texttt{arXiv:1705.07750}
  \item Localizing Moments in Video with Natural Language (2017). \texttt{arXiv:1708.01641}
  \item Dense-Captioning Events in Videos (2017). \texttt{arXiv:1705.00754}
  \item The 2017 DAVIS Challenge on Video Object Segmentation (2017). \texttt{arXiv:1704.00675}
  \item A Style-Based Generator Architecture for Generative Adversarial Networks (2018). \texttt{arXiv:1812.04948}
  \item The Unreasonable Effectiveness of Deep Features as a Perceptual Metric (2018). \texttt{arXiv:1801.03924}
  \item Towards Accurate Generative Models of Video: A New Metric \& Challenges (2018). \texttt{arXiv:1812.01717}
  \item Video-to-Video Synthesis (2018). \texttt{arXiv:1808.06601}
  \item Train Sparsely, Generate Densely: Memory-efficient Unsupervised Training
  of High-resolution Temporal GAN (2018). \texttt{arXiv:1811.09245}
  \item End-to-End Dense Video Captioning with Masked Transformer (2018). \texttt{arXiv:1804.00819}
  \item Generative Modeling by Estimating Gradients of the Data Distribution (2019). \texttt{arXiv:1907.05600}
  \item Exploring the Limits of Transfer Learning with a Unified Text-to-Text
  Transformer (2019). \texttt{arXiv:1910.10683}
  \item HowTo100M: Learning a Text-Video Embedding by Watching Hundred Million
  Narrated Video Clips (2019). \texttt{arXiv:1906.03327}
  \item VideoBERT: A Joint Model for Video and Language Representation Learning (2019). \texttt{arXiv:1904.01766}
  \item Adversarial Video Generation on Complex Datasets (2019). \texttt{arXiv:1907.06571}
  \item VATEX: A Large-Scale, High-Quality Multilingual Dataset for
  Video-and-Language Research (2019). \texttt{arXiv:1904.03493}
  \item Denoising Diffusion Probabilistic Models (2020). \texttt{arXiv:2006.11239}
  \item Denoising Diffusion Implicit Models (2020). \texttt{arXiv:2010.02502}
  \item Score-Based Generative Modeling through Stochastic Differential
  Equations (2020). \texttt{arXiv:2011.13456}
  \item An Image is Worth 16x16 Words: Transformers for Image Recognition at
  Scale (2020). \texttt{arXiv:2010.11929}
  \item Taming Transformers for High-Resolution Image Synthesis (2020). \texttt{arXiv:2012.09841}
  \item Learning Transferable Visual Models From Natural Language Supervision (2021). \texttt{arXiv:2103.00020}
  \item High-Resolution Image Synthesis with Latent Diffusion Models (2021). \texttt{arXiv:2112.10752}
  \item Diffusion Models Beat GANs on Image Synthesis (2021). \texttt{arXiv:2105.05233}
  \item Zero-Shot Text-to-Image Generation (2021). \texttt{arXiv:2102.12092}
  \item Frozen in Time: A Joint Video and Image Encoder for End-to-End Retrieval (2021). \texttt{arXiv:2104.00650}
  \item GLIDE: Towards Photorealistic Image Generation and Editing with
  Text-Guided Diffusion Models (2021). \texttt{arXiv:2112.10741}
  \item Improved Denoising Diffusion Probabilistic Models (2021). \texttt{arXiv:2102.09672}
  \item Cascaded Diffusion Models for High Fidelity Image Generation (2021). \texttt{arXiv:2106.15282}
  \item SDEdit: Guided Image Synthesis and Editing with Stochastic Differential
  Equations (2021). \texttt{arXiv:2108.01073}
  \item VideoGPT: Video Generation using VQ-VAE and Transformers (2021). \texttt{arXiv:2104.10157}
  \item LoRA: Low-Rank Adaptation of Large Language Models (2021). \texttt{arXiv:2106.09685}
  \item N\textbackslash{}"UWA: Visual Synthesis Pre-training for Neural visUal World creAtion (2021). \texttt{arXiv:2111.12417}
  \item GODIVA: Generating Open-DomaIn Videos from nAtural Descriptions (2021). \texttt{arXiv:2104.14806}
  \item A Good Image Generator Is What You Need for High-Resolution Video
  Synthesis (2021). \texttt{arXiv:2104.15069}
  \item StyleGAN-V: A Continuous Video Generator with the Price, Image Quality
  and Perks of StyleGAN2 (2021). \texttt{arXiv:2112.14683}
  \item Hierarchical Text-Conditional Image Generation with CLIP Latents (2022). \texttt{arXiv:2204.06125}
  \item Make-A-Video: Text-to-Video Generation without Text-Video Data (2022). \texttt{arXiv:2209.14792}
  \item Video Diffusion Models (2022). \texttt{arXiv:2204.03458}
  \item Photorealistic Text-to-Image Diffusion Models with Deep Language
  Understanding (2022). \texttt{arXiv:2205.11487}
  \item Classifier-Free Diffusion Guidance (2022). \texttt{arXiv:2207.12598}
  \item Imagen Video: High Definition Video Generation with Diffusion Models (2022). \texttt{arXiv:2210.02303}
  \item CogVideo: Large-scale Pretraining for Text-to-Video Generation via
  Transformers (2022). \texttt{arXiv:2205.15868}
  \item Tune-A-Video: One-Shot Tuning of Image Diffusion Models for
  Text-to-Video Generation (2022). \texttt{arXiv:2212.11565}
  \item Prompt-to-Prompt Image Editing with Cross Attention Control (2022). \texttt{arXiv:2208.01626}
  \item DreamBooth: Fine Tuning Text-to-Image Diffusion Models for
  Subject-Driven Generation (2022). \texttt{arXiv:2208.12242}
  \item MagicVideo: Efficient Video Generation With Latent Diffusion Models (2022). \texttt{arXiv:2211.11018}
  \item Latent Video Diffusion Models for High-Fidelity Long Video Generation (2022). \texttt{arXiv:2211.13221}
  \item Scalable Diffusion Models with Transformers (2022). \texttt{arXiv:2212.09748}
  \item Long Video Generation with Time-Agnostic VQGAN and Time-Sensitive
  Transformer (2022). \texttt{arXiv:2204.03638}
  \item Generating Videos with Dynamics-aware Implicit Generative Adversarial
  Networks (2022). \texttt{arXiv:2202.10571}
  \item LAION-5B: An open large-scale dataset for training next generation
  image-text models (2022). \texttt{arXiv:2210.08402}
  \item An Image is Worth One Word: Personalizing Text-to-Image Generation using
  Textual Inversion (2022). \texttt{arXiv:2208.01618}
  \item Progressive Distillation for Fast Sampling of Diffusion Models (2022). \texttt{arXiv:2202.00512}
  \item MoVideo: Motion-Aware Video Generation with Diffusion Models (2023). \texttt{arXiv:2311.11325}
  \item Hierarchical Spatio-temporal Decoupling for Text-to-Video Generation (2023). \texttt{arXiv:2312.04483}
  \item VideoGen: A Reference-Guided Latent Diffusion Approach for High
  Definition Text-to-Video Generation (2023). \texttt{arXiv:2309.00398}
  \item Stable Video Diffusion: Scaling Latent Video Diffusion Models to Large
  Datasets (2023). \texttt{arXiv:2311.15127}
  \item VideoFusion: Decomposed Diffusion Models for High-Quality Video
  Generation (2023). \texttt{arXiv:2303.08320}
  \item Video Probabilistic Diffusion Models in Projected Latent Space (2023). \texttt{arXiv:2302.07685}
  \item Align your Latents: High-Resolution Video Synthesis with Latent
  Diffusion Models (2023). \texttt{arXiv:2304.08818}
  \item LaMD: Latent Motion Diffusion for Image-Conditional Video Generation (2023). \texttt{arXiv:2304.11603}
  \item Adding Conditional Control to Text-to-Image Diffusion Models (2023). \texttt{arXiv:2302.05543}
  \item VideoChat: Chat-Centric Video Understanding (2023). \texttt{arXiv:2305.06355}
  \item Structure and Content-Guided Video Synthesis with Diffusion Models (2023). \texttt{arXiv:2302.03011}
  \item AnimateDiff: Animate Your Personalized Text-to-Image Diffusion Models
  without Specific Tuning (2023). \texttt{arXiv:2307.04725}
  \item Video-ChatGPT: Towards Detailed Video Understanding via Large Vision and
  Language Models (2023). \texttt{arXiv:2306.05424}
  \item ModelScope Text-to-Video Technical Report (2023). \texttt{arXiv:2308.06571}
  \item Text2Video-Zero: Text-to-Image Diffusion Models are Zero-Shot Video
  Generators (2023). \texttt{arXiv:2303.13439}
  \item FateZero: Fusing Attentions for Zero-shot Text-based Video Editing (2023). \texttt{arXiv:2303.09535}
  \item LAVIE: High-Quality Video Generation with Cascaded Latent Diffusion
  Models (2023). \texttt{arXiv:2309.15103}
  \item VideoComposer: Compositional Video Synthesis with Motion Controllability (2023). \texttt{arXiv:2306.02018}
  \item Video-LLaMA: An Instruction-tuned Audio-Visual Language Model for Video
  Understanding (2023). \texttt{arXiv:2306.02858}
  \item VBench: Comprehensive Benchmark Suite for Video Generative Models (2023). \texttt{arXiv:2311.17982}
  \item Video Diffusion Models: A Survey (2024). \texttt{arXiv:2405.03150}
  \item VideoTetris: Towards Compositional Text-to-Video Generation (2024). \texttt{arXiv:2406.04277}
  \item Grid Diffusion Models for Text-to-Video Generation (2024). \texttt{arXiv:2404.00234}
  \item Redefining Temporal Modeling in Video Diffusion: The Vectorized Timestep
  Approach (2024). \texttt{arXiv:2410.03160}
  \item Decoupled Video Generation with Chain of Training-free Diffusion Model
  Experts (2024). \texttt{arXiv:2408.13423}
  \item Efficient Video Diffusion Models via Content-Frame Motion-Latent
  Decomposition (2024). \texttt{arXiv:2403.14148}
  \item SSM Meets Video Diffusion Models: Efficient Long-Term Video Generation
  with Structured State Spaces (2024). \texttt{arXiv:2403.07711}
  \item Vid3D: Synthesis of Dynamic 3D Scenes using 2D Video Diffusion (2024). \texttt{arXiv:2406.11196}
  \item VideoCrafter2: Overcoming Data Limitations for High-Quality Video
  Diffusion Models (2024). \texttt{arXiv:2401.09047}
  \item Survey of Video Diffusion Models: Foundations, Implementations, and Applications (2025). \texttt{arXiv:2504.16081}
  \item DiffuseSlide: Training-Free High Frame Rate Video Generation Diffusion (2025). \texttt{arXiv:2506.01454}
  \item From Image to Video: An Empirical Study of Diffusion Representations (2025). \texttt{arXiv:2502.07001}
  \item Training-Free Motion-Guided Video Generation with Enhanced Temporal
  Consistency Using Motion Consistency Loss (2025). \texttt{arXiv:2501.07563}
  \item RepVideo: Rethinking Cross-Layer Representation for Video Generation (2025). \texttt{arXiv:2501.08994}
  \item AR-Diffusion: Asynchronous Video Generation with Auto-Regressive
  Diffusion (2025). \texttt{arXiv:2503.07418}
  \item VideoMerge: Towards Training-free Long Video Generation (2025). \texttt{arXiv:2503.09926}
  \item Compositional Video Synthesis by Temporal Object-Centric Learning (2025). \texttt{arXiv:2507.20855}
  \item Temporal In-Context Fine-Tuning with Temporal Reasoning for Versatile Control of Video Diffusion Models (2025). \texttt{arXiv:2506.00996}
  \item Efficient Video Diffusion Models: Advancements and Challenges (2026). \texttt{arXiv:2604.15911}
  \item Image-to-Video Diffusion: From Foundations to Open Frontiers (2026). \texttt{arXiv:2605.17248}
  \item FrameDiT: Diffusion Transformer with Matrix Attention for Efficient Video Generation (2026). \texttt{arXiv:2603.09721}
  \item TeleStyle: Content-Preserving Style Transfer in Images and Videos (2026). \texttt{arXiv:2601.20175}
  \item QwenStyle: Content-Preserving Style Transfer with Qwen-Image-Edit (2026). \texttt{arXiv:2601.06202}
  \item VISTA: Triplet-Supervised Video Style Transfer with Diffusion Transformers (2026). \texttt{arXiv:2605.17312}
\end{enumerate}

\end{document}
